\chapter{Geometry at One Base Point}
\label{ch:geometry}

This chapter answers what geometric data an agent carries at one base point $c_*$ and which of those
data are invariant. The active gauge group at $c_*$ is one copy of $G$ inside the passive vertex group
$\prod_{i\in V}G$ (\Cref{prop:geo-active-image-point}); every independent link carries a
gauge-invariant conjugacy class, and loop holonomy is the telescoped product of these per-edge twists
(\Cref{prop:geo-edge-twist-invariance}); frame agreement, agreement of an equivariant statistic, and
law coherence up to holonomy stabilizers are distinct (\Cref{prop:geo-three-agreements}).
\Cref{ch:generative} uses passive re-expression for covariance of the generative law,
\Cref{ch:selective-communication} builds attention on link-transported comparisons, and
\Cref{ch:network-state} turns holonomy and the trivializing criterion
(\Cref{prop:geo-trivializing-criterion}) into fixed sectors of coarse channels.

\paragraph{Scope.}
Normalizability of a law on links built from the loop penalties of \Cref{rem:geo-loop-penalties} against
the Haar reference is not asserted for a noncompact or general linear $G$ (the compact unitary case is in
the remark, and a Gaussian reference on the Lie algebra normalizes the generator prior of
\Cref{prop:link-generator-prior} for every $G$), and
whether law coherence bears on the selection of a partition into meta-agents is part of
\Cref{open:cg-natural-selector}. \status{NOT-CLAIMED}

\section{The base point, the principal torsor, and associated law fibers}
\label{sec:geo-base}
\label{sec:geo-bundle}
\label{sec:geo-associated}

\definitionheading{Contextual base and base point}{def:geo-context-base}
Let $\mathcal C$ be a finite-dimensional, second countable, Hausdorff smooth manifold whose points
$c$ are called contexts, and fix one base point $c_*\in\mathcal C$.

\definitionheading{Common principal gauge bundle and agent frames}{def:geo-principal-systems}
Let $\pi$ be a smooth principal right $G$-bundle for a Lie group $G$ and, for each agent $i\in V$,
let $u_i$ be one local frame on an open support $\mathcal C_i\ni c_*$:
\begin{gather}
\pi:\mathscr P_G\longrightarrow\mathcal C,
\label{eq:geo-principal-bundle}\\
u_i:\mathcal C_i\longrightarrow\mathscr P_G,
\qquad
\pi\circ u_i=\operatorname{id}_{\mathcal C_i}.
\label{eq:geo-frame-sections}
\end{gather}
The fiber $(\mathscr P_G)_{c_*}$ is a right $G$-torsor and the values $u_i:=u_i(c_*)$ are points of
it. For any reference point $\sigma\in(\mathscr P_G)_{c_*}$ the frames are
$u_i=\sigma\cdot U_i^{-1}$ for unique $U_i\in G$.

The same frame $u_i$ serves the belief and the model channel; \Cref{prop:geo-common-frame-gauge}
shows that this is a gauge choice. On the belief and model sample fibers $(\mathsf K,\mathscr K)$ and
$(\mathsf M,\mathscr M)$ declare law fibers, measurable representations, which need not be equivalent,
faithful, or of equal dimension, and their pushforward actions on laws,
\begin{gather}
\mathcal B_b\subseteq\mathcal P(\mathsf K),
\qquad
\mathcal B_m\subseteq\mathcal P(\mathsf M),
\label{eq:geo-law-fibers}\\
\rho_b:G\longrightarrow\operatorname{Aut}(\mathsf K),
\qquad
\rho_m:G\longrightarrow\operatorname{Aut}(\mathsf M),
\label{eq:geo-representations}\\
\widehat\rho_b(g)q=(\rho_b(g))_\#q,
\qquad
\widehat\rho_m(g)q=(\rho_m(g))_\#q,
\label{eq:geo-pushforward-actions}
\end{gather}
so that $(\widehat\rho_x(g)q)(A)=q(\rho_x(g)^{-1}A)$.

\hypothesisheading{Selected smooth tier}{hyp:geo-smooth-tier}
Whenever a differential statement is made, $\mathcal B_b$ and $\mathcal B_m$ are finite-dimensional
smooth parametrized-measure models, each closed under its $G$-action, and the actions
$(g,q)\mapsto(\rho_x(g))_\#q$\label{eq:geo-smooth-actions} are smooth. The hypothesis excludes
stratified and infinite-dimensional law fibers; a statistical reading of the tier further requires
differentiability in quadratic mean and the integrability each tensor uses, and a Fisher tensor is a
metric only where it is nondegenerate.

\definitionheading{Belief and model associated bundles}{def:geo-associated-bundles}
On invariant law fibers define the bundles below, with projections $\varpi_x:\mathcal E_x\to\mathcal C$
and the quotient convention of the second line,
\begin{gather}
\mathcal E_b=\mathscr P_G\times_{\widehat\rho_b}\mathcal B_b,
\qquad
\mathcal E_m=\mathscr P_G\times_{\widehat\rho_m}\mathcal B_m,
\label{eq:geo-associated-bundles}\\
(u\cdot g,\beta)\sim(u,\widehat\rho_x(g)\beta).
\label{eq:geo-quotient-convention}
\end{gather}
At $c_*$ the frame $u_i$ identifies $\mathcal B_x$ with the fiber $(\mathcal E_x)_{c_*}$ by
$\beta\mapsto[u_i,\beta]$. This is the standard associated-bundle construction
\citep[Section~I.5]{Kobayashi1963}, smooth under \Cref{hyp:geo-smooth-tier}.

\section{Frames, passive re-expression, and the active gauge group at \texorpdfstring{$c_*$}{c*}}
\label{sec:geo-frame-gauge}

A passive re-expression of agent $i$ is a choice $g_i\in G$ that replaces the frame as in the first
line below. Since $[u_i,\beta]=[u_i'\cdot g_i,\beta]=[u_i',\widehat\rho_x(g_i)\beta]$, associated
coordinates change by $\beta'=\widehat\rho_x(g_i)\beta$, and the represented sample-coordinate maps are
those of the second line:
\begin{gather}
u_i'=u_i\cdot g_i^{-1},
\label{eq:geo-local-reframing}\\
R_i^b=\rho_b(g_i),
\qquad
R_i^m=\rho_m(g_i),
\qquad
k_i'=R_i^bk_i,\quad m_i'=R_i^mm_i.
\label{eq:geo-represented-frame-change}
\end{gather}
The group element, not a represented operator, multiplies the frame, and the tuple $(g_i)_{i\in V}$
ranges over the passive vertex group $\prod_{i\in V}G$.

\subsection{Pointwise frame comparisons and the active gauge group}
\label{sec:geo-regime-one}

\propositionheading{Frame comparisons at one base point}{prop:geo-frame-comparisons}
For $i,j\in V$ there is a unique $T_{ij}\in G$ with the first relation below, for every reference
point $\sigma$. The comparisons satisfy $T_{ii}=e_G$, $T_{ji}=T_{ij}^{-1}$, and $T_{ij}T_{jl}=T_{il}$,
and under passive re-expressions they transform by the second line:
\begin{gather}
u_j=u_i\cdot T_{ij},
\qquad
T_{ij}=U_iU_j^{-1},
\label{eq:geo-frame-comparison}\\
T_{ij}'=g_iT_{ij}g_j^{-1},
\qquad
U_i'=g_iU_i.
\label{eq:geo-regime-one-gauge-law}
\end{gather}
\status{ESTABLISHED}

\noindent\emph{Proof.}
Freeness and transitivity of the torsor give uniqueness of $T_{ij}$, and
$u_j=\sigma\cdot U_j^{-1}=u_i\cdot U_iU_j^{-1}$ gives the formula; the identities follow from
uniqueness. From $u_j'=u_j\cdot g_j^{-1}=u_i'\cdot g_iT_{ij}g_j^{-1}$ and
$u_i'=\sigma\cdot U_i^{-1}g_i^{-1}$ follow both laws in \eqref{eq:geo-regime-one-gauge-law}.
\hfill$\square$

\propositionheading{Active gauge image at one base point}{prop:geo-active-image-point}
Let $F\in\operatorname{Aut}_G((\mathscr P_G)_{c_*})$, the group of right-$G$-equivariant bijections of
the principal fiber over $c_*$, have local functions $F(u_i)=u_i\cdot\varphi_i$. Then
\begin{equation}
\varphi_j=T_{ij}^{-1}\varphi_iT_{ij},
\qquad i,j\in V,
\label{eq:geo-active-image-point}
\end{equation}
and $F\mapsto(\varphi_i)_{i\in V}$ is a group isomorphism onto the subgroup
$\mathcal A_{c_*}\subseteq\prod_{i\in V}G$ cut out by \eqref{eq:geo-active-image-point}. For each
$i_0\in V$ the projection $(\varphi_i)_{i\in V}\mapsto\varphi_{i_0}$ is an isomorphism
$\mathcal A_{c_*}\cong G$. In the frame $u_i$, $F$ changes the coordinates of both associated
fibers by $\beta\mapsto\widehat\rho_x(\varphi_i)\beta$, which is the coordinate change of the
passive re-expression $g_i=\varphi_i$. Hence the active transformations at $c_*$ form one copy
$\mathcal A_{c_*}\cong G$ of the passive vertex group, a proper subgroup of $\prod_{i\in V}G$ when
$G$ is nontrivial and $|V|\geq2$, of dimension $\dim G$ against $|V|\dim G$ for a Lie group. The
local functions at $c_*$ of every gauge automorphism of $\mathscr P_G$ form one element of
$\mathcal A_{c_*}$. \status{ESTABLISHED}

\noindent\emph{Proof.}
Equivariance gives $F(u_j)=F(u_i\cdot T_{ij})=u_i\cdot\varphi_iT_{ij}=u_j\cdot T_{ij}^{-1}\varphi_iT_{ij}$,
and uniqueness in the torsor gives \eqref{eq:geo-active-image-point}. Since
$F(F'(u))=F(u\cdot\varphi')=u\cdot\varphi\varphi'$, the map is a homomorphism. The value
$\varphi_{i_0}$ determines $F$ on the whole torsor, so the map is injective, and for
$\varphi_{i_0}\in G$ the map $u_{i_0}\cdot g\mapsto u_{i_0}\cdot\varphi_{i_0}g$ is an equivariant
bijection, which proves $\mathcal A_{c_*}\cong G$. The coordinate statement is
$F[u_i,\beta]=[u_i\cdot\varphi_i,\beta]=[u_i,\widehat\rho_x(\varphi_i)\beta]$, the passive coordinate
change with $g_i=\varphi_i$. If $G\neq\{e_G\}$ and $i\neq j$, the tuple with $\varphi_i=a\neq e_G$ and
all other entries $e_G$ violates \eqref{eq:geo-active-image-point}, so the subgroup is proper. A gauge
automorphism restricts to an equivariant bijection of each fiber. \hfill$\square$

\subsection{The common frame is a gauge choice}
\label{sec:geo-diagonal}

\propositionheading{Two atlases reduce to one by a passive re-expression}{prop:geo-common-frame-gauge}
Suppose belief and model coordinates were given in two frames $u_i^b,u_i^m$ of the same fiber. Then
there is a unique $h_i\in G$ with the first relation below, and under independent re-expressions
$u_i^{b\prime}=u_i^b\cdot(g_i^b)^{-1}$ and $u_i^{m\prime}=u_i^m\cdot(g_i^m)^{-1}$ it transforms by the
second:
\begin{gather}
u_i^m=u_i^b\cdot h_i,
\label{eq:geo-relative-frame}\\
h_i'=g_i^bh_i(g_i^m)^{-1}.
\label{eq:geo-relative-frame-law}
\end{gather}
In particular the re-expression $g_i^b=e_G$, $g_i^m=h_i$ of the model atlas alone makes the two
frames equal. In two atlases the belief and model comparisons, the local functions
$F(u_i^x)=u_i^x\cdot\varphi_i^x$ of a fiber automorphism $F$, and the representatives of a cross
morphism are related by
\begin{gather}
T_{ij}^m=h_i^{-1}T_{ij}^bh_j,
\label{eq:geo-frame-comparison-relation}\\
\varphi_i^m=h_i^{-1}\varphi_i^bh_i,
\label{eq:geo-diagonal-gauge-functions}\\
\Phi_i^{\mathrm{conv}}=\widehat\rho_m(h_i)\circ\Phi_i^{bm},
\label{eq:geo-map-bundle-frame-conversion}
\end{gather}
where $\Phi_i^{bm}$ is the representative of $\Phi$ with source frame $u_i^b$ and target frame
$u_i^m$, and $\Phi_i^{\mathrm{conv}}$ is its representative with both frames equal to $u_i^b$. All
three relations become identities when $h_i=e_G$. \status{ESTABLISHED}

\noindent\emph{Proof.}
Uniqueness of $h_i$ is the torsor property. Substituting the re-expressions gives
$u_i^{m\prime}=u_i^b\cdot h_i(g_i^m)^{-1}=u_i^{b\prime}\cdot g_i^bh_i(g_i^m)^{-1}$, which is
\eqref{eq:geo-relative-frame-law}. The two expressions $u_j^m=u_j^b\cdot h_j=u_i^b\cdot T_{ij}^bh_j$
and $u_j^m=u_i^m\cdot T_{ij}^m=u_i^b\cdot h_iT_{ij}^m$ give
\eqref{eq:geo-frame-comparison-relation}. Equivariance gives
$F(u_i^m)=F(u_i^b)\cdot h_i=u_i^b\cdot\varphi_i^bh_i=u_i^m\cdot h_i^{-1}\varphi_i^bh_i$. Finally,
$\Phi[u_i^b,\beta]=[u_i^m,\Phi_i^{bm}\beta]=[u_i^b,\widehat\rho_m(h_i)\Phi_i^{bm}\beta]$ by
\eqref{eq:geo-quotient-convention}. \hfill$\square$

No energy, kernel, or evaluator reads $h_i$, so the common frame
\begin{equation}
u_i^b=u_i^m=u_i,
\qquad
h_i=e_G,
\qquad
T_{ij}^b=T_{ij}^m=T_{ij}
\label{eq:geo-diagonal-principal}
\end{equation}
is adopted throughout.

\section{Agents and cross-associated-bundle morphisms}
\label{sec:geo-agents}
\label{sec:geo-cross-morphisms}

\definitionheading{Agent}{def:geo-agent}
Agent $i$ is the section-bearing object $\mathcal A^i$ of the first line below, with the sections of
the second line, and a population sample at $c_*$ is a point $y$ of a product of sample fibers:
\begin{gather}
\mathcal A^i=\bigl(\mathcal C_i;q_i^b,q_i^m,u_i\bigr),
\label{eq:geo-agent-object}\\
q_i^b\in\Gamma(\mathcal C_i,\mathcal E_b|_{\mathcal C_i}),
\qquad
q_i^m\in\Gamma(\mathcal C_i,\mathcal E_m|_{\mathcal C_i}),
\label{eq:geo-agent-sections}\\
y=\bigl((k_1,m_1),\ldots,(k_N,m_N)\bigr)\in\prod_{i\in V}(\mathsf K\times\mathsf M).
\label{eq:geo-population-sample}
\end{gather}
The values of $q_i^b,q_i^m$ at $c_*$ are points of $(\mathcal E_b)_{c_*}$ and $(\mathcal E_m)_{c_*}$, whose
coordinates in the frame $u_i$ are laws in $\mathcal B_b$ and $\mathcal B_m$; the same letters denote
these coordinate laws, which \Cref{def:prob-recognition-marginals} identifies with the local recognition
marginals. Agents need not have equal supports.

\definitionheading{Cross morphisms}{def:geo-cross-morphisms}
Cross morphisms are measurable fiber-preserving maps $\Phi:\mathcal E_b\to\mathcal E_m$ and
$\widetilde\Phi:\mathcal E_m\to\mathcal E_b$, $\varpi_m\circ\Phi=\varpi_b$ and
$\varpi_b\circ\widetilde\Phi=\varpi_m$.\label{eq:geo-cross-channel-map} At $c_*$ they are maps between
the fibers $(\mathcal E_b)_{c_*}$ and $(\mathcal E_m)_{c_*}$. They are not principal-bundle morphisms, gauge automorphisms, same-channel transports, or assumed inverses,
and their existence is additional structure. Under \Cref{hyp:geo-smooth-tier} a differential
statement also assumes them smooth. In the frame $u_i$ their representatives
$\Phi_i:\mathcal B_b\to\mathcal B_m$ and $\widetilde\Phi_i:\mathcal B_m\to\mathcal B_b$ are defined
by $\Phi[u_i,\beta]=[u_i,\Phi_i\beta]$ and $\widetilde\Phi[u_i,\beta]=[u_i,\widetilde\Phi_i\beta]$.

\propositionheading{Representatives and induced cross morphisms}{prop:geo-intertwining-cross-map}
The representatives transform under passive re-expression by the first two lines below, and those of
different agents are related by the third:
\begin{gather}
\Phi_i'=\widehat\rho_m(g_i)\circ\Phi_i\circ\widehat\rho_b(g_i)^{-1},
\label{eq:geo-phi-gauge-law}\\
\widetilde\Phi_i'=\widehat\rho_b(g_i)\circ\widetilde\Phi_i\circ\widehat\rho_m(g_i)^{-1},
\label{eq:geo-tildephi-gauge-law}\\
\Phi_i\circ\widehat\rho_b(T_{ij})=\widehat\rho_m(T_{ij})\circ\Phi_j,
\qquad
\widetilde\Phi_i\circ\widehat\rho_m(T_{ij})=\widehat\rho_b(T_{ij})\circ\widetilde\Phi_j.
\label{eq:geo-cross-map-gluing}
\end{gather}
If maps $\Phi_0:\mathcal B_b\to\mathcal B_m$ and $\widetilde\Phi_0:\mathcal B_m\to\mathcal B_b$
satisfy the first line below for every $g\in G$, then the maps of the second line
\begin{gather}
\Phi_0\circ\widehat\rho_b(g)=\widehat\rho_m(g)\circ\Phi_0,
\qquad
\widetilde\Phi_0\circ\widehat\rho_m(g)=\widehat\rho_b(g)\circ\widetilde\Phi_0,
\label{eq:geo-intertwining}\\
[p,\beta]\longmapsto[p,\Phi_0\beta],
\qquad
[p,\beta]\longmapsto[p,\widetilde\Phi_0\beta]
\label{eq:geo-induced-cross-morphisms}
\end{gather}
are well-defined cross morphisms, whose representatives in every frame are $\Phi_0$ and
$\widetilde\Phi_0$; conversely, $[p,\beta]\mapsto[p,\Phi_0\beta]$ is well defined only if $\Phi_0$
satisfies \eqref{eq:geo-intertwining}. Not every cross morphism is induced in this way, and for a
linear specialization with irreducible, inequivalent representations every equivariant linear map
between the sample fibers vanishes, while cross morphisms still exist. \status{ESTABLISHED}

\noindent\emph{Proof.}
From $[u_i,\beta]=[u_i',\widehat\rho_b(g_i)\beta]$ we get
$\Phi[u_i',\beta']=[u_i,\Phi_i\widehat\rho_b(g_i)^{-1}\beta']=[u_i',\widehat\rho_m(g_i)\Phi_i\widehat\rho_b(g_i)^{-1}\beta']$,
which is \eqref{eq:geo-phi-gauge-law}; interchanging the channels gives
\eqref{eq:geo-tildephi-gauge-law}. Evaluating $\Phi[u_j,\beta]=\Phi[u_i,\widehat\rho_b(T_{ij})\beta]$
in both frames gives \eqref{eq:geo-cross-map-gluing}. The representatives $(p\cdot g,\beta)$ and
$(p,\widehat\rho_b(g)\beta)$ of one class have the same image under
\eqref{eq:geo-induced-cross-morphisms} exactly when $\widehat\rho_m(g)\Phi_0\beta=\Phi_0\widehat\rho_b(g)\beta$,
which proves both directions of the induced statement. For a cross morphism that is not induced, let
$G=SO(2)$ act on the laws on $\R^2$ in both channels by the defining representation, fix
$\sigma\in(\mathscr P_G)_{c_*}$, and put $\Phi[\sigma,\beta]=[\sigma,D_\#\beta]$ with
$D=\operatorname{diag}(2,1)$, which is well defined because every class has a unique representative
with principal component $\sigma$. By \eqref{eq:geo-phi-gauge-law} its representative in the frame
$\sigma\cdot g^{-1}$ is $\widehat\rho_m(g)D_\#\widehat\rho_b(g)^{-1}$, and for the rotation $R_{\pi/2}$
one has $D_\#(R_{\pi/2})_\#\delta_{(1,0)}=\delta_{(0,1)}\neq\delta_{(0,2)}=(R_{\pi/2})_\#D_\#\delta_{(1,0)}$,
so the representative depends on the frame. In the linear specialization, the kernel and image of an
equivariant linear map between irreducible representation spaces are invariant subspaces, so the map is
zero or an isomorphism, and inequivalence excludes an isomorphism. \hfill$\square$

Law-level equivariance is weaker than linear equivariance: an equivariant Markov kernel induces an
equivariant map on laws without supplying a linear map between the sample fibers.

\section{Independent links at \texorpdfstring{$c_*$}{c*}: vertex law, holonomy, edge twists}
\label{sec:geo-regime-two}
\label{sec:geo-trivializing}
\label{sec:phil-agreements}

\definitionheading{Typed graph links}{def:geo-graph-links}
Let $\mathfrak I=(V,E,\bar{\ })$ be a finite interaction multigraph at $c_*$: a finite set $E$ of
oriented edge copies $e:j\to i$ with a fixed-point-free reversal $e\mapsto\bar e$, parallel copies
allowed. For each channel $x\in\{b,m\}$ and each $e:j\to i$ declare a link $\Theta_e^x$ as in the
first line below, with link map $\mathsf t_e^x(u_j\cdot g)=u_i\cdot\Theta_e^xg$ of the fiber; endpoint
notation $\Theta_{ij}^x$ is used only when the copy $j\to i$ is unique. Under passive re-expression the
links obey the vertex law of the second and third lines, with $R_i^x=\rho_x(g_i)$, because
$\mathsf t_e^x(u_j'\cdot g)=u_i\cdot\Theta_e^xg_j^{-1}g=u_i'\cdot g_i\Theta_e^xg_j^{-1}g$. For a walk
$p=(e_0,\ldots,e_{r-1})$ from $j$ to $i$, that is, with $e_a:i_{a+1}\to i_a$, $i_0=i$, and $i_r=j$, the
path product and, when $i=j$, the holonomy are those of the fourth line, and the common-link
specialization of the fifth line declares that the two link maps coincide in the common frame:
\begin{gather}
\Theta_e^x\in G,
\qquad
\Theta_{\bar e}^x=(\Theta_e^x)^{-1},
\label{eq:geo-regime-two-links}\\
\Theta_e^{x\prime}=g_i\Theta_e^xg_j^{-1},
\label{eq:geo-regime-two-gauge-law}\\
\rho_x(\Theta_e^{x\prime})=R_i^x\rho_x(\Theta_e^x)(R_j^x)^{-1},
\label{eq:gen-represented-gauge-links}\\
\Theta_p^x=\Theta_{e_0}^x\Theta_{e_1}^x\cdots\Theta_{e_{r-1}}^x,
\qquad
\mathrm{Hol}^x(p)=\Theta_p^x,
\label{eq:geo-link-holonomy}\\
\Theta_e^b=\Theta_e^m=\Theta_e.
\label{eq:geo-common-link}
\end{gather}

\propositionheading{Per-edge twist invariance at one base point}{prop:geo-edge-twist-invariance}
Fix a reference point $\sigma$ with $u_i=\sigma\cdot U_i^{-1}$, and define the per-edge twist
$\mathring\Theta_e^x$ of $e:j\to i$ by the first line below; a closed walk $p$ based at $i_0$ then has
the holonomy of the second line:
\begin{gather}
\mathsf t_e^x(\sigma)=\sigma\cdot\mathring\Theta_e^x,
\qquad
\mathring\Theta_e^x=U_i^{-1}\Theta_e^xU_j,
\label{eq:geo-edge-twist}\\
\mathrm{Hol}^x(p)=U_{i_0}\mathring\Theta_{e_0}^x\mathring\Theta_{e_1}^x\cdots\mathring\Theta_{e_{r-1}}^xU_{i_0}^{-1},
\qquad
\mathrm{Hol}^x(p)'=g_{i_0}\mathrm{Hol}^x(p)g_{i_0}^{-1}.
\label{eq:geo-edge-twist-holonomy}
\end{gather}
Then $\mathsf t_e^x$ is a right-$G$-equivariant bijection that does not depend on the frames, and
$\mathsf t_{\bar e}^x=(\mathsf t_e^x)^{-1}$. Every twist is unchanged by arbitrary passive
re-expressions; a change of reference point $\sigma'=\sigma\cdot g_0$ conjugates every twist by the
same $g_0$; and an active $F$ with $F(\sigma)=\sigma\cdot\kappa$ sends $\mathsf t_e^x$ to
$F\circ\mathsf t_e^x\circ F^{-1}$ and every twist to $\kappa\mathring\Theta_e^x\kappa^{-1}$. Each edge
copy therefore carries the conjugacy class $[\mathring\Theta_e^x]$, invariant under all passive and
active changes, and two link data sets on the same multigraph are related by one active
automorphism exactly when their twist tuples, computed from one reference point, are related by one
simultaneous conjugation. The conjugacy class of $\mathrm{Hol}^x(p)$ and the condition
$\mathrm{Hol}^x(p)=e_G$ are invariant. \status{ESTABLISHED}

\noindent\emph{Proof.}
Every point of the fiber is $u_j\cdot g$ for a unique $g$, so $\mathsf t_e^x$ is a well-defined
equivariant bijection whose inverse $u_i\cdot g\mapsto u_j\cdot(\Theta_e^x)^{-1}g$ is $\mathsf t_{\bar e}^x$,
and by the computation in \Cref{def:geo-graph-links} re-expressed frames and links define the same
map. Since $\sigma=u_j\cdot U_j$, $\mathsf t_e^x(\sigma)=u_i\cdot\Theta_e^xU_j=\sigma\cdot
U_i^{-1}\Theta_e^xU_j$; the left side depends only on $\mathsf t_e^x$ and $\sigma$, which proves passive
invariance, and $\mathsf t_e^x(\sigma\cdot g_0)=(\sigma\cdot g_0)\cdot g_0^{-1}\mathring\Theta_e^xg_0$.
For the active statement,
$F\mathsf t_e^xF^{-1}(\sigma)=F(\sigma\cdot\mathring\Theta_e^x\kappa^{-1})=\sigma\cdot\kappa\mathring\Theta_e^x\kappa^{-1}$;
an equivariant self-map of the torsor is determined by its value at $\sigma$, and $F\mapsto\kappa$ is a
bijection onto $G$, which gives the criterion for two link data sets. Finally
$\Theta_{e_a}^x=U_{i_a}\mathring\Theta_{e_a}^xU_{i_{a+1}}^{-1}$ telescopes, and $U_{i_0}'=g_{i_0}U_{i_0}$ by
\eqref{eq:geo-regime-one-gauge-law}. \hfill$\square$

\propositionheading{Trivializing criterion}{prop:geo-trivializing-criterion}
\label{prop:gauss-coboundary-trivialization}
Fix a subfamily $I\subseteq V$ and elements $\Xi_e\in G$ on the internal oriented edge copies of $I$.
The following are equivalent:
\begin{equation}
\begin{gathered}
\exists V_i\in G\ (i\in I)
\quad
\Xi_e=V_iV_j^{-1}
\quad\text{for every internal }e:j\to i,\\
\Xi_{e_0}\Xi_{e_1}\cdots\Xi_{e_{r-1}}=e_G
\quad\text{for every closed internal walk }(e_0,\ldots,e_{r-1}).
\end{gathered}
\label{eq:geo-trivializing-criterion}
\end{equation}
When they hold, the potentials $V_i$ are unique up to right multiplication by one element of $G$ on
each connected component of the internal multigraph. The statement and proof use only the group
structure, so they hold for any group in place of $G$, in particular for a subgroup of $\GL(K)$
containing represented links. \status{ESTABLISHED}

\noindent\emph{Proof.}
The first condition telescopes around every closed walk. Conversely, backtracks
$(e,\bar e)$ are closed walks, so $\Xi_{\bar e}=\Xi_e^{-1}$ and reversed walks have inverse products.
Choose a base vertex in each connected component and set $V_i=\Xi_{p}$ for an internal walk $p$ from
that vertex to $i$, with products ordered as in \eqref{eq:geo-link-holonomy}; two choices differ by a
closed walk and agree, and extending a walk to $j$ by $e:j\to i$ gives $V_i=\Xi_eV_j$. If also
$\Xi_e=W_iW_j^{-1}$, then $W_i^{-1}V_i=W_j^{-1}V_j$ along every internal edge, so $W_i^{-1}V_i$ is
constant on each component. \hfill$\square$

\definitionheading{Graph-flat, jointly trivializing, and frame-aligned links}{def:geo-graph-flat-links}
The links of channel $x$ are graph flat on $I\subseteq V$ when \Cref{prop:geo-trivializing-criterion}
holds for $\Xi_e=\Theta_e^x$, and $I$ is jointly trivializing when this holds in both channels. The
links are frame aligned\label{hyp:geo-flat-links} when $\Theta_e^x=T_{ij}=U_iU_j^{-1}$ for every
$e:j\to i$,\label{eq:geo-flat-regime} equivalently when every twist $\mathring\Theta_e^x$ equals $e_G$
and every link map $\mathsf t_e^x$ is the identity of the fiber.

\remarkheading{Frame alignment, graph flatness, and per-edge classes}{rem:geo-flatness-strictness}
Frame alignment implies graph flatness, because $T_{ij}=U_iU_j^{-1}$ is a coboundary, but not
conversely: on two vertices joined by one edge copy $e:2\to1$ and its reversal, every closed walk is
a backtrack, so any $\Theta_e^x\neq T_{12}$ is graph flat with potentials $V_1=\Theta_e^x$,
$V_2=e_G$, while $[\mathring\Theta_e^x]\neq\{e_G\}$. Graph flatness of the links is graph flatness
of the twists, since $\Theta_e^x=V_iV_j^{-1}$ exactly when $\mathring\Theta_e^x=W_iW_j^{-1}$ with
$W_i=U_i^{-1}V_i$. A single link coordinate between distinct vertices carries no invariant: the passive re-expression
$g_i=e_G$, $g_j=\Theta_e^x$ sends $\Theta_e^x$ to $e_G$. On a tree every loop product is trivial,
yet the per-edge classes can all differ from $\{e_G\}$. Graph flatness is a property of the links
alone: the one fiber over $c_*$ is always a trivial torsor, links on a triangle with nontrivial loop
product are not graph flat, and all-identity links are graph flat whatever the topology of
$\mathscr P_G$. \status{ESTABLISHED}

\remarkheading{Invariant loop penalties}{rem:geo-loop-penalties}
Let $p$ be a closed walk based at $i_0$ and $\mathsf U=\rho_x(\mathrm{Hol}^x(p))$. If $\rho_x$ is unitary
of dimension $n$, then
$\frac{1}{2n}\operatorname{Re}\Tr(2I_n-\mathsf U-\mathsf U^*)=\frac{1}{2n}\Tr((I_n-\mathsf U)^*(I_n-\mathsf U))\geq0$,
with equality exactly when $\mathsf U=I_n$, and the value is invariant under passive re-expression,
which conjugates $\mathsf U$ by the unitary $R_{i_0}^x=\rho_x(g_{i_0})$. For $n\geq2$ no norm on
$\mathfrak{gl}(n,\R)$ is invariant under conjugation by all of $\GL(n,\R)$: conjugation by
$\operatorname{diag}(2,1,\ldots,1)$ doubles the matrix unit $E_{12}$. A positive invariant penalty for
general linear links therefore needs further data: if a symmetric positive definite fiber metric
$\mathsf G^{\mathrm{fib}}$ at $i_0$ is declared to transform by
$\mathsf G^{\mathrm{fib}}\mapsto(R_{i_0}^x)^{-\top}\mathsf G^{\mathrm{fib}}(R_{i_0}^x)^{-1}$, then
$\Tr\bigl((\mathsf G^{\mathrm{fib}})^{-1}(I_n-\mathsf U)^\top\mathsf G^{\mathrm{fib}}(I_n-\mathsf U)\bigr)$
is invariant, nonnegative, and zero exactly when $\mathsf U=I_n$. For compact $G$ with $\rho_x$
unitary, the penalty $\frac{1}{2n}\Tr((I_n-\mathsf U)^*(I_n-\mathsf U))$ of the first sentence lies
in $[0,2]$, so for finitely many closed walks $p$ with holonomies
$\mathsf U_p=\rho_x(\mathrm{Hol}^x(p))$ and coefficients $\mathrm c_p\geq0$ the law on links with
density $\exp\bigl(-\sum_p\mathrm c_p\frac{1}{2n}\Tr((I_n-\mathsf U_p)^*(I_n-\mathsf U_p))\bigr)$
against the product, over one orientation of each edge copy, of the Haar probability measure of $G$
has normalizer in $[e^{-2\sum_p\mathrm c_p},1]$. \status{ESTABLISHED}

\noindent\emph{Proof.}
Unitarity gives $(I_n-\mathsf U)^*(I_n-\mathsf U)=2I_n-\mathsf U-\mathsf U^*$, whose trace is real and
equals the squared Frobenius norm of $I_n-\mathsf U$. With $\Delta=I_n-\mathsf U$ and
$\mathsf G=\mathsf G^{\mathrm{fib}}$, passive re-expression sends $\Delta$ to
$R_{i_0}^x\Delta(R_{i_0}^x)^{-1}$, cyclicity of the trace gives invariance, and
$\Tr(\mathsf G^{-1}\Delta^\top\mathsf G\Delta)=\lVert\mathsf G^{1/2}\Delta\mathsf G^{-1/2}\rVert_{\mathrm F}^2$.
A norm $\lVert\cdot\rVert_\star$ with $\lVert E_{12}\rVert_\star>0$ cannot satisfy
$\lVert2E_{12}\rVert_\star=\lVert E_{12}\rVert_\star$. For the compact case,
$\Tr(2I_n-\mathsf U-\mathsf U^*)=2n-2\operatorname{Re}\Tr\mathsf U$ and $\lvert\Tr\mathsf U\rvert\leq n$
for unitary $\mathsf U$, so that penalty lies in $[0,2]$ and the exponent in $[-2\sum_p\mathrm c_p,0]$;
each $\mathsf U_p$ is a continuous function of the links, and integrating against a probability
measure gives the two bounds on the normalizer. \hfill$\square$

In the exponential chart of \Cref{ch:links}, with the fiber metric at each vertex induced by one matrix at
the reference point, one half of the general linear penalty above is the loop defect of
\Cref{def:link-loop-defect} (\Cref{prop:link-loop-defect}); and for $G=\GL(K,\R)$ or $\GL^+(K)$, on a
closed walk that traverses some reversal pair exactly once, every continuous penalty of the holonomy that
is unchanged by passive re-expression takes at every unipotent holonomy its value at trivial holonomy
(\Cref{cor:link-no-invariant-defect}), so a penalty that vanishes only at trivial holonomy needs data
beyond the holonomy, such as the fiber metric.

\definitionheading{Law coherence}{def:geo-law-coherence}
Let $I\subseteq V$ induce a connected internal multigraph. Laws $q_i\in\mathcal B_x$, $i\in I$, are law
coherent on $I$ when
\begin{equation}
q_i=\widehat\rho_x(\Theta_e^x)q_j
\qquad\text{on each internal edge copy }e:j\to i.
\label{eq:phil-law-coherence}
\end{equation}
For a $G$-set $S$ and a map $s:\mathcal B_x\to S$ with $s(\widehat\rho_x(g)q)=g\cdot s(q)$, the laws are
$s$-agreeing when $s(q_i)=\Theta_e^x\cdot s(q_j)$ on each internal edge copy; the mean of laws on a
linear representation space is the basic instance. The links have frame agreement on $I$ when they
are graph flat on $I$.

\propositionheading{Three agreements}{prop:geo-three-agreements}
Let $I$ be as in \Cref{def:geo-law-coherence} and fix $i_0\in I$. (i) The laws are law coherent
exactly when $q_{i_0}$ is fixed by $\widehat\rho_x(\mathrm{Hol}^x(p))$ for every closed internal walk
$p$ based at $i_0$ and $q_i=\widehat\rho_x(\Theta^x_{p_i})q_{i_0}$ for some, hence every, internal walk
$p_i$ from $i_0$ to $i$; for an equivariant statistic the same holds with $s(q_{i_0})$ fixed by the
holonomy, so that frame agreement and mean agreement on edges, the two conditions of the first line
below, are conditions on the links and on the means, and mean agreement holds exactly when the base
mean satisfies the second line and the other means are its transports:
\begin{gather}
\mathrm{Hol}^x(p)=e_G\ \text{for every closed internal walk }p,
\quad
\mu_i=\rho_x(\Theta_e^x)\mu_j\ \text{on each internal }e:j\to i,
\label{eq:phil-two-agreements}\\
\mu_{i_0}\in\bigcap_{p\text{ closed at }i_0}\ker\bigl(\rho_x(\mathrm{Hol}^x(p))-\operatorname{id}\bigr).
\label{eq:phil-agreement-kernel}
\end{gather}
(ii) Law coherence implies $s$-agreement for every equivariant statistic $s$, and for the mean the converse fails. (iii) Frame agreement neither implies
nor is implied by law coherence or by mean agreement. (iv) With several independent cycles, trivial
holonomy or a fixed base mean for one loop is necessary but not sufficient. In particular, for
$G=\GL(2,\R)$ acting on laws on $\R^2$, three agents with the common law $\mathcal N(0,\Sigma)$,
$\Sigma=\begin{psmallmatrix}4&1\\1&1\end{psmallmatrix}$, and links
$\Theta_{12}=\Sigma^{1/2}R(0.7)\Sigma^{-1/2}$ and $\Theta_{23}=\Theta_{31}=I_2$, with $R(\phi)$ the
rotation by $\phi$, are law coherent while the triangle holonomy $\Theta_{12}$ is not the identity, is
not orthogonal, and fixes no nonzero vector. \status{ESTABLISHED}

\noindent\emph{Proof.}
(i) Composing \eqref{eq:phil-law-coherence} along an internal walk $p$ from $j$ to $i$ gives
$q_i=\widehat\rho_x(\Theta^x_p)q_j$, which yields both conditions. Conversely, let the
conditions hold; for $e:j\to i$ and walks $p_i,p_j$ from $i_0$, the walk that follows $p_j$, then $e$,
then $p_i$ backward is closed at $i_0$ with holonomy $(\Theta_{p_i}^x)^{-1}\Theta_e^x\Theta_{p_j}^x$, so
$\widehat\rho_x(\Theta_e^x)q_j=\widehat\rho_x(\Theta^x_{p_i})\widehat\rho_x\bigl((\Theta_{p_i}^x)^{-1}\Theta_e^x\Theta_{p_j}^x\bigr)q_{i_0}=q_i$.
The statistic version is the same argument in $S$, and the vectors left unchanged by the linear map
$\rho_x(\mathrm{Hol}^x(p))$ form $\ker(\rho_x(\mathrm{Hol}^x(p))-\operatorname{id})$. (ii) Apply $s$ to
\eqref{eq:phil-law-coherence}. For the converse take two agents joined by one edge copy with identity
link and the laws $\mathcal N(0,I_2)$ and $\mathcal N(0,2I_2)$: the means agree and the laws do not.
(iii) The same two agents with the laws $\mathcal N(0,I_2)$ and $\mathcal N(v,I_2)$, $v\neq0$, have frame
agreement without mean agreement or law coherence. In the displayed witness,
$\Theta_{12}\Sigma\Theta_{12}^\top=\Sigma^{1/2}R(0.7)R(0.7)^\top\Sigma^{1/2}=\Sigma$, so every transported
law equals $\mathcal N(0,\Sigma)$ and all means vanish, while $\Theta_{12}$ is similar to a rotation by
$0.7$, hence has no eigenvalue $1$, and $\Theta_{12}\Theta_{12}^\top\neq I_2$ because $\Sigma$ is not
a multiple of $I_2$. (iv) Two vertices joined by three parallel copies with $\Theta_e=\Theta_f\neq\Theta_g$
have trivial holonomy on the walk through $e$ and $\bar f$ and nontrivial holonomy on the walk through
$e$ and $\bar g$. \hfill$\square$

\section{Running example: three agents on one loop}
\label{sec:geo-running-example}

Take $G=SO(3)$, $\rho_b$ the defining representation on $\mathsf K=\R^3$, and $\rho_m$ the trivial
representation on $\mathsf M=\R$, so the channels carry inequivalent representations. Let
$V=\{1,2,3\}$, let $\mathfrak I$ be the triangle with edge copies $e_1:2\to1$, $e_2:3\to2$, $e_3:1\to3$
and their reversals, and, with $R_z(\phi)$ and $R_x(\phi)$ the rotations by $\phi$ about the third and
first coordinate axes, declare the belief links of the first line below; model links act trivially.
Choose $\sigma=u_1$ and, for display, $u_2=u_3=u_1$, so $U_i=e_G$ and each twist equals its link. The
per-edge classes are the angles $\pi/2$, $\pi/2$, and $0$, and the loop $p=(e_1,e_2,e_3)$ based at
agent $1$ has the holonomy of the second line,
\begin{gather}
\Theta_{e_1}^b=R_z(\tfrac\pi2),
\qquad
\Theta_{e_2}^b=R_x(\tfrac\pi2),
\qquad
\Theta_{e_3}^b=I_3,
\label{eq:geo-running-links}\\
\mathrm{Hol}^b(p)=R_z(\tfrac\pi2)R_x(\tfrac\pi2),
\qquad
\Tr\mathrm{Hol}^b(p)=0,
\label{eq:geo-running-holonomy}
\end{gather}
a rotation by $2\pi/3$; two quarter turns about one axis would give a rotation by $\pi$ with the same
classes, so the class tuple does not determine the holonomy, while the tuple modulo one simultaneous
conjugation does (\Cref{prop:geo-edge-twist-invariance}). A passive re-expression $g_2$ changes
$\Theta_{e_1}^b$ to $\Theta_{e_1}^bg_2^{-1}$ and $\Theta_{e_2}^b$ to $g_2\Theta_{e_2}^b$ but changes no
twist, recomputed with $U_2'=g_2$, and not the holonomy. The loop is not graph flat, and the isotropic
laws $q_i^b=\mathcal N(0,s^2I_3)$ are law coherent on it, because every rotation fixes them
(\Cref{prop:geo-three-agreements}). One record reads the transported discrepancy of agents $1$ and $2$:
the Markov kernel from $\mathsf K\times\mathsf K$ to $\{0,1\}$ with
\begin{equation}
K_{a_1}\bigl(k_1,k_2,\{1\}\bigr)=\exp\Bigl(-\tfrac12\bigl\lvert k_1-\rho_b(\Theta_{e_1}^b)k_2\bigr\rvert^2\Bigr),
\label{eq:geo-running-record}
\end{equation}
which is gauge inert, since passive re-expressions rotate $k_1-\rho_b(\Theta_{e_1}^b)k_2$ by
$\rho_b(g_1)$.

\section*{Physicist's summary}

At one base point all frames lie in one principal fiber, so they are compared directly by
$T_{ij}=U_iU_j^{-1}$, the only active symmetry is one global rotation of that fiber, and per-agent frame
changes are bookkeeping. Each independent link carries an invariant, its conjugacy class after dressing
by the frames, whereas in lattice gauge theory only closed-loop products are invariant
\citep{WilsonConfinement1974}. Agreement comes in three kinds, flat links, agreement of a statistic such
as the mean, and law coherence, which survives nontrivial holonomy when the base law has a nontrivial
stabilizer.
